\color{unidarkgreen}{\subsection{Контроль синхронизма и напряжений}}
\color{black}

\begin{enumerate}[label=\arabic{section}.\arabic{subsection}.\arabic*,  align=left, labelsep=4pt, leftmargin=0pt, itemindent=57pt]

\item
Функциональный блок контроля синхронизма и напряжений (КСН) предназначен для формирования
разрешающих сигналов, соответствующих условиям включения для АПВ и оперативного включения.

\begin{figure}[H]
\centering
\includegraphics[width=0.8\textwidth,height=0.8\textheight,keepaspectratio]{\fbpath/761.3351 КСН/img/KSN.pdf}
\captionof{figure}{Функциональная схема КСН В1(ЛВ) ВН}\label{pic:KSN}
\end{figure}

\item
Ввод в работу КСН осуществляется с помощью программного ключа \textbf{«КСН»}. Режимы включения с
контролем синхронизма (КС) и с улавливанием синхронизма (УС) вводятся соответствующими программными
ключами \textbf{«КС»} и \textbf{«УС»}.

\item
Если разность частот напряжений 1СШ и 2СШ не превышает уставку \textbf{«df макс КС»}, то разрешается
включение с контролем синхронизма. Для режима включения с КС также контролируются следующие условия:
\begin{itemize}
\item наличие напряжения на обеих системах шин;
\item разность напряжений на системах шин, не превышающая по модулю значения заданной уставки;
\item разность фаз напряжений шин и линии, не превышающая по модулю значения заданной уставки.
\end{itemize}

\item
При выводе программного ключа \textbf{«КС»} включение разрешается без контроля частоты, а только по
уровню напряжений и разности их фаз.

\item
Если разность частот превышает уставку \textbf{«df макс КС»}, но не выходит за пределы уставки \textbf{«df пред.
доп.»}, то разрешается включение с улавливанием синхронизма. Сигнал включения с улавливанием
синхронизма формируется с опережением на время включения выключателя. Режим включения с УС можно
вывести программным ключом \textbf{«УС»}.

\item
Наличие или отсутствие напряжения на первой системе шин контролируется измерительными органами
максимального и минимального напряжения, пороги срабатывания которых задаются уставками \textbf{«Uнал Ш1
КСН»} и \textbf{«Uотс Ш1 КСН»} соответственно. Пороги срабатывания измерительных органов для контроля
напряжения на второй системе шин задаются уставками \textbf{«Uнал Ш2 КСН»} и \textbf{«Uотс Ш2 КСН»}.

\item
Предусмотрена блокировка разрешающих сигналов КСН при не симметрии в цепях напряжения, для
1СШ величина не симметрии задается с помощью уставок \textbf{«U2макс Ш1 КСН»} и \textbf{«3U0макс Ш1 КСН»}, для 2СШ
с помощью уставок \textbf{«U2макс Ш2 КСН»} и \textbf{«3U0макс Ш2 КСН»}.

\item
Разрешение оперативного включения выключателя с контролем синхронизма производится с помощью команды управления \textit{«Опер. Вкл. В с КС(УС)»}.

\item
Оперативный ввод/вывод КСН производится с помощью БУ (блоков управления функциями) \textit{«Вывод КСН»}. 

\item
Основные режимы включения и контролируемые для них условия приведены в таблице \ref{ksn:tbl1}.


\small
\begin{longtable}{|>{\centering\arraybackslash}m{5.3cm}|>{\centering\arraybackslash}m{11.6cm}|}
% Первая страница: полный заголовок
\caption{Условия включения\hfill\vspace{-0.5\baselineskip}}\label{ksn:tbl1}\\
\hline
\rowcolor{gray!30}
Режим включения & Условия включения \\
\hline
\endfirsthead

% Промежуточные страницы: «Продолжение таблицы»
\caption*{\hspace{3pt}\emph{Продолжение таблицы \ref{ksn:tbl1}}\hfill\ }\vspace{-0.5\baselineskip} \\
\hline
\rowcolor{gray!30}
Режим включения & Условия включения \\
\hline
\endhead

ОНШ1+ННШ2 &   
			\begin{itemize}
			\vspace{3pt}
			\item $ \textrm{U}_\textrm{Ш1} $ меньше уставки \textbf{«Uотс Ш1 КСН»};
			\item $ \textrm{U}_\textrm{Ш2} $ больше уставки \textbf{«Uнал Ш2 КСН»}.
			\vspace{3pt}
			\end{itemize} \\
\hline
ОНШ2+ННШ1 & 
			\begin{itemize}
			\vspace{3pt}
			\item $ \textrm{U}_\textrm{Ш1} $ больше уставки \textbf{«Uнал Ш1 КСН»};
			\item $ \textrm{U}_\textrm{Ш2} $ меньше уставки \textbf{«Uотс Ш2 КСН»}.
			\vspace{3pt}
			\end{itemize} \\
\hline
КС(УС)	&
			\vspace{3pt}
			
			\begin{itemize}
			\item[] \noindent 1) Включение контролем синхронизма:
			\item $ \textrm{U}_\textrm{Ш1} $ больше уставки \textbf{«Uнал Ш1 КСН»};
			\item $ \textrm{U}_\textrm{Ш2} $ больше уставки \textbf{«Uнал Ш2 КСН»};
			\item разность напряжений 1СШ и 2СШ не превышает по модулю значения заданной уставки \textbf{«dU макс КС(УС)»};
			\item разность фаз напряжений 1СШ и 2СШ не превышает по модулю значения заданной уставки \textbf{«dФ макс КС»};
			\item разность частот напряжений 1СШ и 2СШ не превышает значения заданной уставки \textbf{«df макс КС»};
			\end{itemize}
			
			\begin{itemize}
			\item[] \noindent  2) Включение с улавливанием синхронизма:
			\item $ \textrm{U}_\textrm{Ш1} $ больше уставки \textbf{«Uнал Ш1 КСН»};
			\item $ \textrm{U}_\textrm{Ш2} $ больше уставки \textbf{«Uнал Ш2 КСН»};
			\item разность напряжений 1СШ и 2СШ не превышает по модулю значения заданной уставки \textbf{«dU макс КС(УС)»};
			\item разность частот напряжений 1СШ и 2СШ превышает значение уставки \textbf{«df макс КС»};
			\item разность частот напряжений 1СШ и 2СШ не превышает значения заданной уставки \textbf{«df пред доп»};
			\item за время \textbf{«Т опереж вкл В»} с момента подачи сигнала на включение выключателя расхождение углов между векторами напряжений на 1СШ и 2СШ не превысит 10º.
			\vspace{3pt}
			\end{itemize} \\
\hline

\end{longtable}
\normalsize	

\item
Сигнал разрешения АПВ 1СШ \textit{«Разр. АПВ 1СШ»} формируется при выполнении условий включения в
режимах ОНШ1+ННШ2, КС или УС.

\item
Сигнал разрешения АПВ 2СШ \textit{«Разр. АПВ 2СШ»} формируется при выполнении условий включения в
режимах ОНШ2+ННШ1, КС или УС. 

\item
Сигнал разрешения оперативного включения \textit{«Разр. опер. вкл. В с КС»} может формироваться от
внутренней логики КСН или от внешнего сигнала в зависимости от положения программного переключателя
\textbf{«Контр. опер. вкл. В»}.

\item
В положении \textit{«Внутр.»} оперативное включение с КС разрешается при выполнении условий включения в
режимах ОНШ1+ННШ2 (при опробовании 1СШ), ОНШ2+ННШ1 (при опробовании 2СШ), КС или УС (при
замыкании в транзит). Также предусмотрена возможность оперативного включения с контролем отсутствия
напряжения на 1СШ или на 2СШ при введенном программном ключе \textbf{«Опер. вкл. В с КОН»}. При отсутствии
внешнего сигнала \textit{«Опер. вкл. с КС(УС)»} оперативное включение может выполняться без контроля напряжений
и условий синхронизма.

\item
В положении \textit{«Внеш.»} оперативное включение разрешается при наличии внешнего сигнала 
\textit{«Внеш. сигнал КС(УС)»}.

\small
\begin{longtable}{|>{\centering\arraybackslash}m{5.3cm}|>{\centering\arraybackslash}m{3.3cm}|>{\centering\arraybackslash}m{4.2cm}|>{\centering\arraybackslash}m{1.8cm}|>{\centering\arraybackslash}m{1cm}|}
\caption{Перечень уставочных параметров КСН \hfill\vspace{-0.5\baselineskip}}\label{ksn:tbl2}\\ 
\hline
\rowcolor{gray!30}
Параметр & Условное обозначение на схеме & Значение/ Диапазон & Единица измерения & Шаг \\ 
\hline
\endfirsthead
\caption*{\hspace{3pt}\emph{Продолжение таблицы \ref{ksn:tbl2}\hfill\vspace{-0.5\baselineskip}}} \\ % сделано по ГОСТ 2.105 п.6.8.7
\hline
\rowcolor{gray!30}
Параметр & Условное обозначение на схеме & Значение/ Диапазон & Единица измерения & Шаг \\ 
%\hline
\endhead

%===m>КСН|-
%===m
\end{longtable}
\normalsize

\end{enumerate}